\chapter{Combinaison de signaux}\label{ch:rs:signal-combination}

Une équipe de recherche dispose de quarante signaux, chacun doté d'un \omterm{def:rs:anatomy-of-a-predictor:ic}{coefficient d'information} proche de 0.015. Combinés à poids égaux, ils atteignent 0.064.
Pondérés par leurs IC passés, ils atteignent 0.099 sur un marché simulé et 0.075 sur un autre~; avec les poids de manuel qui maximisent le
ratio d'information du mélange, 0.064 et 0.041. Dans un univers de cent valeurs avec un an d'historique, les poids de régression ne rapportent presque rien dans l'un des deux
marchés, et là, rien ne bat les poids égaux. Les prévisionnistes connaissent ce phénomène sous le nom d'énigme~: Stock et Watson (2004) ont établi que les combinaisons les plus
performantes de prévisions de croissance de la production, comme la moyenne, étaient les moins sensibles aux performances récentes des
prévisions individuelles. Ce chapitre construit les mélanges (\texttt{firm.combine}), explique ce que suppose chacun et mesure sur des données simulées quand les
poids estimés valent leur erreur d'estimation.

\section{La géométrie d'un mélange linéaire}

\begin{definition}[Combinaison de signaux, signal composite, mélange équipondéré]\label{def:rs:signal-combination:blend}
La \emph{combinaison de signaux}\index{combinaison de signaux} est la construction d'un prédicteur à partir de plusieurs. Un \emph{signal
composite}\index{signal composite} est une somme pondérée $F = \sum_k w_k z_k$ des signaux $z_k$, chacun standardisé sur l'ensemble des valeurs à chaque date~;
le \emph{mélange équipondéré}\index{mélange équipondéré} fixe chaque $w_k = 1/K$.
\end{definition}

\begin{proposition}[L'IC d'un mélange équipondéré]\label{prop:rs:signal-combination:geometry}
Si $K$ signaux standardisés ont chacun une corrélation $\rho$ avec la cible et une corrélation deux à deux $c$ entre eux, leur \omterm{def:rs:signal-combination:blend}{mélange équipondéré}
a une corrélation $K\rho/\sqrt{K + K(K - 1)c}$ avec la cible, qui croît avec $K$ vers $\rho/\sqrt c$.
\end{proposition}

\begin{proof}
La covariance de $\sum_k z_k$ avec la cible (standardisée) vaut $K\rho$~; sa variance vaut $K + K(K-1)c$. On divise, puis on fait tendre $K \to \infty$.
\end{proof}

Le plafond est fixé par la corrélation entre signaux, et non par leur nombre~: quarante signaux qui se répètent forment en réalité un seul signal. Le
premier marché du chapitre (\texttt{rs\_combine}, le monde \emph{stable}) injecte la structure~: chaque cible mensuelle est le rendement du mois suivant
d'une valeur sur \texttt{firm.synthmkt} (corrigé du marché, standardisé sur les 704 valeurs cotées sur toute la période) plus un rendement espéré issu de huit
thèmes latents, dont cinq portent de l'information (de 0.10 à 0.03) et trois aucune~; chacun des quarante signaux mesure un thème, cinq
signaux par thème, à travers un bruit que partage son thème et un bruit propre dont l'amplitude croît de 1 à 4 au sein du thème. Sur les cinq
dernières années, l'IC moyen des signaux vaut 0.015 et leur corrélation moyenne deux à deux 0.027. La combinaison mesurée de $K$ signaux pris dans un ordre aléatoire suit
de près la proposition (la \cref{fig:rs:signal-combination:geometry})~: 0.021 pour deux, 0.042 pour dix, 0.064 pour les quarante contre 0.065 donné
par la formule, en tendant vers un plafond de 0.090.

\begin{omfigure}
\begin{tikzpicture}
\begin{axis}[width=10cm,height=5.4cm,xlabel={nombre de signaux combinés, $K$},ylabel={IC hors échantillon},tick label style={font=\small},
  label style={font=\small},xmin=0,xmax=41,ymin=0,ymax=0.1,grid=major,grid style={gray!25},
  yticklabel style={/pgf/number format/fixed,/pgf/number format/precision=2},scaled y ticks=false,legend pos=south east,
  legend style={font=\footnotesize},table/col sep=comma]
\addplot[omDef,very thick,mark=*,mark size=1.6pt] table[x=k,y=measured]{figdata/research/14-signal-combination/geometry.csv};
\addplot[omThm,thick,dashed] table[x=k,y=theory]{figdata/research/14-signal-combination/geometry.csv};
\addplot[black!50,thin,dotted] table[x=k,y=ceiling]{figdata/research/14-signal-combination/geometry.csv};
\legend{mesure (100 ordres aléatoires),$K\bar\rho/\sqrt{K + K(K-1)\bar c}$,plafond $\bar\rho/\sqrt{\bar c}$}
\end{axis}
\end{tikzpicture}

\omcaption{IC hors échantillon du \omterm{def:rs:signal-combination:blend}{mélange équipondéré} de $K$ des quarante signaux injectés (monde stable, cinq dernières années), comparé à la
proposition, avec leur IC moyen de 0.015 et leur corrélation moyenne deux à deux de 0.027. Données~: \texttt{rs\_combine} sur
\texttt{firm.synthmkt}.}\label{fig:rs:signal-combination:geometry}
\end{omfigure}

\section{Poids issus des estimations}

\begin{definition}[Mélange pondéré par l'IC, mélange à ICIR maximal]\label{def:rs:signal-combination:weights}
Un \emph{mélange pondéré par l'IC}\index{mélange pondéré par l'IC} fixe chaque poids proportionnel à l'IC moyen passé du signal (nul s'il est négatif). Un
\emph{mélange à ICIR maximal}\index{mélange à ICIR maximal} fixe $w \propto \Sigma^{-1}\mu$, avec $\mu$ le vecteur des IC moyens passés des signaux et
$\Sigma$ la matrice de covariance de leurs séries temporelles d'IC~: les poids qui maximisent la moyenne de l'IC du mélange rapportée à son écart type.
\end{definition}

Les poids égaux ne supposent rien des signaux~; les poids pondérés par l'IC supposent que leur qualité diffère et persiste~; la règle à ICIR maximal suppose, en
outre, que la matrice de covariance de quarante séries d'IC, estimée à partir de 59 mois, est assez précise pour être inversée. Le second marché (le
monde \emph{en dérive}) supprime les deux premières hypothèses~: chaque signal a le même bruit, et l'information de chaque thème oscille autour
moyenne commune de 0.04, selon une autorégression de demi-vie d'environ deux ans, si bien que les poids égaux sont justes en moyenne et que tout poids estimé
poursuit le bruit et une information périmée. Ajustés sur les cinq premières années et évalués sur les cinq dernières, avec les 704 valeurs~:

\begin{center}
\small
\setlength{\tabcolsep}{4pt}%
\begin{tabular}{@{}lcccc@{}}
\toprule
 & \multicolumn{2}{c}{monde stable} & \multicolumn{2}{c}{monde en dérive}\\
mélange & en échantillon & hors échantillon & en échantillon & hors échantillon\\
\midrule
poids égaux & 0.065 & 0.064 & 0.055 & 0.060\\
poids pondérés par l'IC & 0.091 & 0.099 & 0.073 & 0.075\\
ICIR maximal, sans rétrécissement & 0.050 & 0.064 & 0.052 & 0.041\\
ICIR maximal, $\Sigma$ diagonale & 0.090 & 0.098 & 0.074 & 0.074\\
régression groupée (MCO) & 0.093 & 0.100 & 0.076 & 0.071\\
lasso & 0.093 & 0.102 & 0.075 & 0.071\\
la vérité terrain (poids oracle) & --- & 0.102 & --- & 0.087\\
\bottomrule
\end{tabular}
\end{center}

Avec 704 valeurs par mois, les poids pondérés par l'IC captent l'essentiel de ce qu'offre la vérité terrain dans le monde stable, et battent encore les poids égaux dans le
monde en dérive (une information qui dérive pendant deux ans reste présente un an plus tard). La règle à ICIR maximal non rétrécie est le pire mélange estimé dans les deux cas~:
elle inverse une matrice de covariance de quarante séries bruitées et charge ses plus petites valeurs propres. En réduisant cette matrice à sa diagonale, c'est-à-dire celui du livre~4 (chapitres 14 et 22) poussé à sa limite, on reste à environ un millième du résultat des poids pondérés par l'IC dans les deux mondes (0.098 contre 0.099, 0.074 contre 0.075).

\section{Combinaison régularisée}

Les poids d'une régression groupée de la cible sur les quarante signaux sont estimés à partir de $59 \times 704$ observations, et là, la
régression est aussi bonne que les poids pondérés par l'IC. L'énigme apparaît lorsque l'échantillon se réduit. Dans un univers de 100 valeurs avec douze mois d'historique
(la \cref{fig:rs:signal-combination:shrink}), les poids de régression non pénalisés rapportent un IC hors échantillon de 0.073 dans le monde stable et de 0.005 dans
monde en dérive, contre 0.069 pour les poids égaux dans les deux cas. Rétrécir les poids de régression vers les poids égaux,
$w_\lambda = (1 - \lambda)\hat w + \lambda/K$, est la régularisation la plus simple, et $\lambda$ admet un optimum~: 0.4 dans le monde stable (IC 0.081),
1 dans le monde en dérive (poids égaux). Smith et Wallis (2009) ont localisé l'énigme précisément là~: l'erreur d'échantillon fini dans l'estimation des
poids de combinaison, qui peut dépasser le gain que procure leur estimation.

\begin{omfigure}
\begin{tikzpicture}
\begin{axis}[width=10cm,height=5.6cm,xlabel={rétrécissement vers les poids égaux, $\lambda$},ylabel={IC hors échantillon},tick label style={font=\small},
  label style={font=\small},xmin=0,xmax=1,ymin=0,ymax=0.1,grid=major,grid style={gray!25},
  yticklabel style={/pgf/number format/fixed,/pgf/number format/precision=2},scaled y ticks=false,legend columns=2,
  legend style={font=\footnotesize,at={(0.5,-0.25)},anchor=north,draw=none,fill=none,
  /tikz/every even column/.append style={column sep=8pt}},table/col sep=comma]
\addplot[omDef,very thick,mark=*,mark size=1.5pt] table[x=lam,y=stable]{figdata/research/14-signal-combination/shrink.csv};
\addplot[omThm,very thick,mark=square*,mark size=1.5pt] table[x=lam,y=drifting]{figdata/research/14-signal-combination/shrink.csv};
\addplot[omDef,thin,dashed] table[x=lam,y=stable_icw]{figdata/research/14-signal-combination/shrink.csv};
\addplot[omThm,thin,dashed] table[x=lam,y=drifting_icw]{figdata/research/14-signal-combination/shrink.csv};
\legend{monde stable,monde en dérive,poids pondérés par l'IC (stable),poids pondérés par l'IC (en dérive)}
\end{axis}
\end{tikzpicture}

\omcaption{Un univers de 100 titres, mélanges calibrés sur douze mois et jugés sur les soixante suivants~: IC hors échantillon de la régression groupée
dont les poids sont ramenés vers des poids égaux ($\lambda = 1$), et des poids d'IC (en pointillés). L'optimum du monde stable est intérieur ($\lambda = 0.4$)~; celui
du monde mouvant correspond à des poids égaux. Données~: \texttt{rs\_combine.small\_universe}.}\label{fig:rs:signal-combination:shrink}
\end{omfigure}

\begin{definition}[Énigme de la combinaison des prévisions]\label{def:rs:signal-combination:puzzle}
L'\emph{énigme de la combinaison des prévisions}\index{énigme de la combinaison des prévisions} est le constat empirique répété selon lequel des combinaisons simples de
prévisions (la moyenne) surpassent les combinaisons dont les poids sont estimés.
\end{definition}

La fenêtre compte moins que le monde~: avec 6, 12, 24 et 48 mois d'historique, les poids d'IC du monde stable battent les poids égaux (0.085, 0.086,
0.082 et 0.092 contre 0.069), tandis que ceux du monde mouvant perdent jusqu'à quatre ans d'historique (0.043, 0.049, 0.067, puis 0.074). Ridge et lasso
pénalisent la régression vers zéro plutôt que vers des poids égaux~; sur un panel standardisé, cela réduit l'échelle du composite davantage que sa
direction, et sur le panel complet elles ont modifié l'IC hors échantillon d'au plus 0.003 (Ridge avec une pénalité de 1 dans le monde mouvant). La pénalité qui compte est celle qui tend vers le bon a priori,
et pour des signaux de qualité voisine et une information instable, le bon a priori consiste en des poids égaux.

\section{Stacking}

\begin{definition}[Stacking]\label{def:rs:signal-combination:stacking}
Le \emph{stacking}\index{stacking} combine plusieurs prédicteurs (ici, plusieurs mélanges), avec des poids estimés sur leurs prévisions pour des données non utilisées pour
les construire~: les régressions empilées de Breiman (1996) utilisent des prévisions validées par validation croisée et des moindres carrés sous contraintes de
non-négativité, idée qu'il attribue à Wolpert (1992).
\end{definition}

Des plis ordonnés dans le temps sont indispensables~: un mélange calibré sur des mois postérieurs à ceux qu'il prédit constitue une fuite. \texttt{firm.stack} calibre chaque
mélange de base (poids égaux, poids d'IC, poids de régression) sur une fenêtre croissante, recueille ses prévisions pour le bloc de mois suivant, et
régresse la cible sur ces prévisions sans poids négatifs. Sur le panel complet, le stack a placé tout son poids sur les poids d'IC dans les deux mondes, ce qui correspond
à leur IC hors échantillon (0.099 et 0.075)~: le \omterm{def:rs:signal-combination:stacking}{stacking} choisit entre les mélanges, et il ne vaut que le meilleur d'entre eux et la longueur de ses
plis.

\section{Orthogonalisation}

\begin{definition}[Orthogonalisation, orthogonalisation symétrique]\label{def:rs:signal-combination:orth}
L'\emph{orthogonalisation}\index{orthogonalisation} transforme un ensemble de signaux, à chaque date, en signaux non corrélés engendrant le même espace~:
\emph{séquentiellement} (Gram--Schmidt), chaque signal étant résidualisé sur ceux qui le précèdent, de sorte que l'ordre détermine celui qui conserve la part commune~; ou
par \emph{orthogonalisation symétrique}\index{orthogonalisation symétrique}, $Z S^{-1/2}$ avec $S$ la matrice de corrélation des signaux, l'ensemble orthonormé
le plus proche des signaux d'origine, une construction qui remonte à l'orthonormalisation de Löwdin des orbitales atomiques qui se recouvrent (1950).
\end{definition}

Les signaux orthogonalisés rendent les poids lisibles (chaque poids est la contribution marginale d'une information qu'aucun autre signal ne porte) et rendent
séparables les expositions du modèle de risque (résidualiser sur les expositions aux facteurs d'un portefeuille revient à la même opération). Ils ne créent pas d'information. Des
poids égaux sur les signaux orthogonalisés symétriquement, recalculés chaque mois, donnent un IC hors échantillon de 0.064 dans le monde stable, soit le même que
les signaux bruts (0.064), chaque signal orthogonalisé étant corrélé à 0.95 avec son original~; l'\omterm{def:rs:signal-combination:orth}{orthogonalisation} séquentielle dans l'ordre des
ICs dans l'échantillon donne 0.050, car le meilleur signal conserve la part commune de son thème et les autres sont réduits à leur bruit avant
d'être pondérés à égalité avec lui.

\section{Fiches de prédicteur}

\begin{predictorcard}{Composite de quarante signaux}\label{pred:rs:signal-combination:composite}
\sfield{Définition} $\sum_k w_k z_k$ avec $z_k$ standardisé par date et $w$ les poids de régression ramenés vers des poids égaux par $\lambda$ choisi sur
des plis ordonnés dans le temps.
\sfield{Entrées et horodatage} Chaque signal tel qu'il est connu à la date de décision~; les poids recalibrés uniquement sur les dates qui le précèdent.
\sfield{Justification} Le plafond des poids égaux $\bar\rho/\sqrt{\bar c}$ (\cref{prop:rs:signal-combination:geometry})~; les poids ne se justifient que lorsque la
qualité diffère de façon persistante.
\sfield{Horizon et demi-vie} Ceux des composantes, pondérés.
\sfield{Normalisation} Standardiser le composite par date~; neutraliser au portefeuille (chapitre~6).
\sfield{Modes de défaillance} Poids d'inverse de covariance non réduits~; recalibrage avec anticipation (look-ahead)~; un composite dominé par les signaux répétés d'un même thème.
\sfield{Sources} Stock et Watson (2004)~; Smith et Wallis (2009)~; \texttt{rs\_combine}.
\end{predictorcard}

\section{Tutoriel~: quarante signaux faibles}\label{tut:rs:signal-combination:tut}

\begin{tutorial}
\textbf{Objectif.} Combiner quarante signaux simulés dans deux mondes en appliquant toutes les règles du chapitre, dans l'échantillon et hors échantillon, sur le panel complet et sur un petit
univers, et trouver la réduction vers des poids égaux qui l'emporte. \textbf{État final~:}
les \cref{fig:rs:signal-combination:geometry,fig:rs:signal-combination:shrink}~; le tableau.
\begin{tutsteps}
\item \textbf{Poids à ICIR maximal}, avec la covariance réduite vers sa diagonale.
\omcode{code/firm/combine/firm_combine.py}{68}{75}{Poids à ICIR maximal avec une covariance réduite des IC.}\label{lst:rs:signal-combination:icir}
\item \textbf{\omterm{def:rs:signal-combination:orth}{Orthogonalisation symétrique}} de la section transversale d'une date.
\omcode{code/firm/combine/firm_combine.py}{157}{161}{Les signaux orthonormés les plus proches des originaux.}\label{lst:rs:signal-combination:sym}
\item \textbf{Exécuter} \texttt{results()} pour les deux mondes, \texttt{small\_universe()} pour des fenêtres de 6 à 48 mois, \texttt{orthogonalised()},
\texttt{geometry()} et \texttt{fig\_combine.py}.
\end{tutsteps}
\textbf{Pour aller plus loin.} Faire passer la demi-vie du monde mouvant à six mois et observer les poids d'IC perdre même sur le panel complet~; donner à trois
signaux un IC négatif et comparer les poids d'IC (qui les écartent) aux poids égaux (qui les conservent).
\end{tutorial}

\section{Construction~: le mélangeur}\label{bld:rs:signal-combination:combine}

\begin{build}
\bfield{Rôle} Un lieu unique où les signaux de la société deviennent un composite, avec toutes les règles du chapitre derrière la même interface, et la
comparaison hors échantillon qui permet de choisir entre eux.
\bfield{Interface} \texttt{zscore}, \texttt{ic\_series}, \texttt{ic\_matrix}, \texttt{equal}, \texttt{blend(X, w)}, \texttt{ic\_weights},
\texttt{max\_icir(ic, shrink)}, \texttt{toward\_equal(w, lam)}, \texttt{ridge}, \texttt{lasso}, \texttt{nnls}, \texttt{stack(preds, y, folds)},
\texttt{time\_folds}, \texttt{orth\_sequential}, \texttt{orth\_symmetric}, \texttt{residualise\_on}.
\bfield{Règles} Les signaux sont standardisés par date avant l'application de tout poids~; les poids et les hyperparamètres sont estimés sur les dates qui précèdent celles sur lesquelles
ils sont jugés~; le \omterm{def:rs:signal-combination:blend}{mélange équipondéré} est toujours présenté en parallèle.
\bfield{Tests d'acceptation} \texttt{code/firm/combine/tests/}~: z-scores et ICs calculés à la main~; un panel simulé où le bon signal obtient le plus grand poids d'IC,
le \omterm{def:rs:signal-combination:weights}{mélange à ICIR maximal} diagonal égale la moyenne divisée par la variance, Ridge réduit, lasso annule le signal de bruit~; NNLS calculé à la main~; un stack qui retient la
base informative~; les deux \omterm{def:rs:signal-combination:orth}{orthogonalisations} sont orthonormées, la symétrique est la plus proche des originaux, la séquentielle conserve sa première colonne.
\bfield{Pour aller plus loin} Poids selon le régime~; mélanges hiérarchiques (au sein des thèmes, puis entre eux)~; poids bayésiens avec un a priori aux poids égaux.
\end{build}

\begin{omsources}
\item J.~H.~Stock et M.~W.~Watson, «~Combination forecasts of output growth in a seven-country data set~», \emph{Journal of Forecasting}
23(6), 2004.
\item J.~Smith et K.~F.~Wallis, «~A simple explanation of the forecast combination puzzle~», \emph{Oxford Bulletin of Economics and Statistics}
71(3), 2009.
\item L.~Breiman, «~Stacked regressions~», \emph{Machine Learning} 24, 1996~; D.~H.~Wolpert, «~Stacked generalization~», \emph{Neural Networks}
5(2), 1992.
\item P.-O.~Löwdin, «~On the non-orthogonality problem connected with the use of atomic wave functions in the theory of molecules and
crystals~», \emph{Journal of Chemical Physics} 18(3), 1950.
\end{omsources}

\section{Exercices}

\begin{exercise}[$\star$]\label{exo:rs:signal-combination:1}
Dix signaux ont chacun un IC de 0.02 et une corrélation deux à deux de 0.25. Quel est l'IC de leur \omterm{def:rs:signal-combination:blend}{mélange équipondéré}, et sa limite lorsque le nombre de
tels signaux croît~?
\end{exercise}

\begin{exercise}[$\star$]\label{exo:rs:signal-combination:2}
Trois signaux ont pour IC moyens passés 0.03, 0.01 et $-0.01$. Donner leurs pondérations par l'IC, ainsi que leurs pondérations rétrécies de moitié vers les pondérations équipondérées.
\end{exercise}

\begin{exercise}[$\star$]\label{exo:rs:signal-combination:3}
Pourquoi les plis du \omterm{def:rs:signal-combination:stacking}{stacking} doivent-ils être ordonnés dans le temps pour des données financières, alors que le \omterm{def:rs:signal-combination:stacking}{stacking} de Breiman utilisait une validation croisée ordinaire~?
\end{exercise}

\begin{exercise}[$\star\star$]\label{exo:rs:signal-combination:4}
Deux signaux ont pour IC $\rho_1 = 0.03$, $\rho_2 = 0.02$ et pour corrélation $c = 0.5$. Déterminer les pondérations du meilleur mélange linéaire (pour la corrélation avec la
cible) et son IC, et comparer avec les pondérations équipondérées.
\end{exercise}

\begin{exercise}[$\star\star$]\label{exo:rs:signal-combination:5}
Montrer que l'\omterm{def:rs:signal-combination:orth}{orthogonalisation} séquentielle de deux signaux standardisés de corrélation $c$ laisse le premier inchangé et remplace le second par
$(z_2 - cz_1)/\sqrt{1 - c^2}$. Quel est l'IC du nouveau second signal en fonction de $\rho_1$, $\rho_2$ et $c$~?
\end{exercise}

\begin{exercise}[$\star\star$]\label{exo:rs:signal-combination:6}
Pourquoi l'inversion de la matrice de covariance de quarante séries d'IC estimées sur soixante mois produit-elle des pondérations extrêmes~? Relier cela aux valeurs propres
d'une matrice de covariance empirique (livre~4, chapitre~22).
\end{exercise}

\begin{exercise}[$\star\star\star$]\label{exo:rs:signal-combination:7}
\emph{Programmation.} Avec \texttt{small\_universe}, déterminer le $\lambda$ optimal dans le monde stable pour des fenêtres de 6, 24 et 48 mois. L'optimum se rapproche-t-il
des pondérations de régression à mesure que la fenêtre s'allonge~?
\end{exercise}

\begin{exercise}[$\star\star\star$]\label{exo:rs:signal-combination:8}
\emph{Trouver la faille.} «~Nous avons ajusté les pondérations du mélange par régression sur les dix années complètes, et elles ont dépassé les pondérations équipondérées de 40~\% sur la même période.~»
\end{exercise}

\section{Problème~: Quarante signaux faibles}

\begin{problem}[{problème du week-end --- quand les pondérations valent la peine d'être estimées}]\label{pb:rs:signal-combination:1}
Les deux mondes simulés du chapitre sur \texttt{firm.synthmkt}~: quarante signaux répartis en huit thèmes, 704 valeurs, dix ans de données mensuelles scindées en deux~; puis
un univers de 100 valeurs.

\medskip\noindent\textbf{Partie I --- La géométrie.}
\begin{enumerate}
\item Quels sont l'IC moyen hors échantillon des signaux et leur corrélation moyenne deux à deux dans le monde stable~?
\item Que donne la \cref{prop:rs:signal-combination:geometry} pour quarante signaux, et qu'a-t-on mesuré~?
\item Quel est le plafond, et pourquoi n'est-il pas atteint en ajoutant des signaux~?
\item Quelle caractéristique volontairement intégrée rend les signaux corrélés~?
\end{enumerate}

\medskip\noindent\textbf{Partie II --- Panel complet.}
\begin{enumerate}[resume]
\item Donner les IC hors échantillon des pondérations équipondérées, des pondérations par l'IC et de la vérité de référence dans les deux mondes.
\item Pourquoi la règle à ICIR maximal, sans rétrécissement, se comporte-t-elle le plus mal~?
\item Que donne le rétrécissement de sa matrice de covariance vers la diagonale~?
\item Pourquoi les pondérations par l'IC battent-elles encore les pondérations équipondérées dans le monde à dérive~?
\end{enumerate}

\medskip\noindent\textbf{Partie III --- Petit univers.}
\begin{enumerate}[resume]
\item Avec 100 valeurs et douze mois, donner l'IC hors échantillon des pondérations de régression non pénalisées dans chaque monde.
\item Énoncer le \emph{résultat nommé}~: l'IC hors échantillon du meilleur mélange face aux pondérations équipondérées, et le rétrécissement vers les pondérations équipondérées qui
le maximise, dans chaque monde.
\item Comment les pondérations par l'IC se comparent-elles aux pondérations équipondérées sur 6, 12, 24 et 48 mois dans chaque monde~?
\item Quelle explication de l'énigme Smith et Wallis ont-ils identifiée~?
\end{enumerate}

\medskip\noindent\textbf{Partie IV --- Au-delà.}
\begin{enumerate}[resume]
\item Quelles pondérations le \omterm{def:rs:signal-combination:stacking}{stacking} a-t-il retenues, et pourquoi ne pouvait-il pas faire mieux~?
\item Que donnent les pondérations égales sur des signaux orthogonalisés de façon symétrique et de façon séquentielle, et pourquoi diffèrent-elles~?
\item Quand une équipe doit-elle orthogonaliser ses signaux~?
\item Comment choisirait-on $\lambda$ en production~?
\item Qu'est-ce qui ferait des pondérations égales un mauvais a priori~?
\item Pourquoi les pénalisations ridge et lasso modifient-elles si peu le résultat sur panel complet ici~?
\item Que devrait publier chaque mélange à côté de son propre IC~?
\item En une phrase~: quand les pondérations estimées valent-elles la peine d'être estimées~?
\end{enumerate}
\end{problem}

\section{Questions d'entretien}

\begin{interviewq}[$\star$ \iqroles{chercheur}]\label{iq:rs:signal-combination:1}
Vous disposez de cinquante signaux alpha. Comment les combinez-vous~?
\end{interviewq}

\begin{interviewq}[$\star\star$ \iqroles{chercheur}]\label{iq:rs:signal-combination:2}
Qu'est-ce que l'\omterm{def:rs:signal-combination:puzzle}{énigme de la combinaison des prévisions}, et qu'est-ce qui l'explique~?
\end{interviewq}

\begin{interviewq}[$\star\star$ \iqroles{chercheur, ingénieur ML}]\label{iq:rs:signal-combination:3}
Comment utiliseriez-vous le \omterm{def:rs:signal-combination:stacking}{stacking} sur des signaux alpha sans fuite de données~?
\end{interviewq}

\begin{interviewq}[$\star\star$ \iqroles{chercheur}]\label{iq:rs:signal-combination:4}
L'ajout d'un nouveau signal d'IC 0.02 n'a pas amélioré votre \omterm{def:rs:signal-combination:blend}{signal composite}. Pourquoi cela pourrait-il être le cas~?
\end{interviewq}

\begin{interviewq}[$\star\star$ \iqroles{chercheur, risque}]\label{iq:rs:signal-combination:5}
Quelle est la différence entre \omterm{def:rs:signal-combination:orth}{orthogonalisation} séquentielle et \omterm{def:rs:signal-combination:orth}{orthogonalisation symétrique}, et quand utiliseriez-vous chacune~?
\end{interviewq}

\begin{interviewq}[$\star\star\star$ \iqroles{chercheur}]\label{iq:rs:signal-combination:6}
Dériver les pondérations qui maximisent la corrélation d'un mélange linéaire avec la cible, étant donnés les IC des signaux et leur matrice de corrélation, et expliquer
pourquoi leur estimation est dangereuse.
\end{interviewq}
